\documentclass[aspectratio=169, 10pt]{beamer}
\usepackage{beamer_theme_cienciadados2026}

\title[Aula 08 - I/O de Arquivos]{Ciência dos Dados I: Análise Exploratória}
\subtitle{Aula 08: Entrada/Saída de Arquivos e Leitura de Dados Compactados}
\author{Prof. Dr. Ney Lemke}
\institute{UNESP -- Instituto de Biociências de Botucatu}
\date{Versão 2026}

\begin{document}

\frame{\titlepage}

\begin{frame}{Sumário da Aula}
  \tableofcontents
\end{frame}

\section{Gerenciadores de Contexto}

\begin{frame}[fragile]{Manipulação Segura com \texttt{with open}}
  Nunca abra arquivos sem o gerenciador de contexto:
  \begin{lstlisting}
# Modo seguro: garante o fechamento do arquivo mesmo se ocorrer um erro
with open('dados.txt', 'r', encoding='utf-8') as f:
    for linha in f:
        print(linha.strip())
  \end{lstlisting}
  \begin{block}{Por que especificar \texttt{encoding='utf-8'}?}
    Sistemas operacionais diferentes adotam codificações distintas por padrão (Windows costumava adotar cp1252; Linux adota UTF-8). Definir UTF-8 explicitamente garante portabilidade.
  \end{block}
\end{frame}

\section{Formatos de Dados: Texto, CSV e Compactação}

\begin{frame}[fragile]{Lendo Arquivos Compactados sem Descompactar no Disco}
  Muitas séries temporais e dados genômicos são armazenados em formato comprimido (\texttt{.gz}) para economizar espaço de armazenamento:
  \begin{lstlisting}
import gzip

with gzip.open('arq_abril_outono.txt.gz', 'rt', encoding='utf-8') as f:
    primeira_linha = f.readline()
    print(primeira_linha)
  \end{lstlisting}
  \begin{exampleblock}{Vantagem em Ciência de Dados}
    A descompressão ocorre em tempo real diretamente na memória durante a leitura de streaming, economizando dezenas de gigabytes de espaço em disco.
  \end{exampleblock}
\end{frame}

\section{O Módulo Nativo \texttt{csv}}

\begin{frame}[fragile]{Leitura Estruturada com \texttt{csv.reader} e \texttt{DictReader}}
  \begin{lstlisting}
import csv

with open('matriznum.csv', 'r', encoding='utf-8') as f:
    leitor = csv.reader(f)
    for linha in leitor:
        # linha e uma lista de strings: ['1.5', '2.3', '4.1']
        valores = [float(x) for x in linha if x]
  \end{lstlisting}
\end{frame}

\begin{frame}{Resumo da Aula}
  \begin{itemize}
    \item Gerenciadores de contexto evitam vazamentos de descritores de arquivo no sistema operacional.
    \item Arquivos compactados podem ser lidos diretamente sem extração prévia.
    \item \textbf{Prática}: Leitura e sumarização de arquivos de texto, CSVs e dados compactados com gzip.
  \end{itemize}
\end{frame}

\end{document}
